\section{Related Work}
\label{Sec2:RelatedWork}

\subsection{Autoregressive and Scale-wise Visual Generation}
Autoregressive (AR) modeling, fundamentally rooted in the probability chain rule, has long served as a cornerstone of generative AI by decomposing intractable high-dimensional distributions into sequences of conditional probabilities. Early pixel-level paradigms, exemplified by PixelRNN~\cite{van2016pixel} and PixelCNN++~\cite{salimans2017pixelcnn++}, theoretically validated this approach but were practically constrained by prohibitive $O(N)$ sequential complexity and limited receptive fields over flattened sequences. The field witnessed a paradigm shift with the introduction of Vector Quantized Variational Autoencoders (VQ-VAE)~\cite{van2017neural} and VQGAN~\cite{esser2021taming}. By compressing high-resolution visual signals into low-dimensional discrete codebooks, these methods enabled Transformers~\cite{vaswani2017attention} to model long-range dependencies within a computationally feasible latent space. This ``Discrete Tokenization + Transformer'' recipe has since become the gold standard for large-scale synthesis, powering foundational models such as DALL-E~\cite{ramesh2021zero}, VideoGPT~\cite{yan2021videogpt}, and MAGVIT~\cite{yu2023language}. Despite the scalability of latent AR models, the conventional raster-scan ordering (top-left to bottom-right) inherently imposes a unidirectional bias that neglects the 2D spatial locality of visual data. To address this structural limitation, recent works like Visual Autoregressive modeling (VAR)~\cite{tian2024visual} and our preliminary framework SAMPO~\cite{wang2025sampo} pioneered a next-scale prediction paradigm. By reformulating generation as a hierarchical progression from coarse structural maps to fine-grained details, these approaches not only align with the global-to-local nature of human perception but also drastically compress sequence complexity from $O(n^2)$ to $O(\log n)$. Nevertheless, a critical vulnerability persists across this lineage: the reliance on discrete quantization. The discrete codebook acts as a lossy information bottleneck, introducing irreducible quantization artifacts and precluding the precise modeling of continuous physical dynamics, such as fluid turbulence or subtle illumination gradients. SAMPO++ fundamentally transcends this limitation by replacing the discrete tokenizer with a unified continuous flow formulation, thereby enabling infinite-resolution synthesis while retaining the efficient structural planning of scale-wise autoregression.


\subsection{Flow Matching and Continuous Latent Dynamics}
Continuous generative models operate on the principle of transforming a simple prior distribution (\textit{e.g.}, Gaussian) into a complex data distribution. Diffusion Probabilistic Models (DPMs)~\cite{ho2020denoising, song2020score} achieve this by learning to reverse a gradual noising process. While DPMs and Latent Diffusion Models (LDM)~\cite{rombach2022high} have established the benchmark for perceptual quality (\textit{e.g.}, Sora~\cite{liu2024sora}), their stochastic sampling trajectories are typically curved and complex, necessitating extensive discretization steps for high-quality inference. Flow Matching (FM)~\cite{lipman2023flow, liuflow, albergo2025stochastic,wang2025flowram} has emerged as a robust alternative, grounded in the theory of Continuous Normalizing Flows (CNFs). Unlike diffusion, FM objectives facilitate the construction of Conditional Vector Fields (CVFs) that regress linear Optimal Transport (OT) paths between source and target distributions. This formulation yields straighter generation trajectories, significantly faster convergence, and superior stability for ODE solvers. Recent variants, such as Rectified Flow~\cite{liuflow} and SiT~\cite{ma2024sit}, have further optimized this framework utilizing Scalable Interpolant Transformers.

An emerging research direction involves hybridizing AR architectures with continuous objectives to reconcile semantic reasoning with visual generation. Approaches like Transfusion~\cite{zhou2024transfusion} and Show-o~\cite{xie2024show} demonstrate the feasibility of training Transformers on mixed modalities, such as discrete text and continuous image patches. Most relevant to our work is FlowAR~\cite{ren2024flowar}, which applies flow matching to the residual of scale-wise image generation. However, FlowAR remains confined to static imagery, functioning primarily as an image upsampler rather than a dynamic simulator. To address this limitation, SAMPO++ extends hybrid continuous-discrete paradigm into the spatiotemporal domain. By integrating temporal autoregression, our framework explicitly enforces the strict causality essential for world modeling, moving beyond the synthesis of independent high-fidelity frames.


\subsection{Action-Conditioned World Models for Embodied Agents}
The World Model paradigm~\cite{ha2018world} posits that embodied agents utilize internal predictive simulators of the environment to facilitate planning and policy learning. Early RSSM-based instantiations, such as the Dreamer framework~\cite{hafner2019dream, hafner2020dreamerv2, hafner2023dreamerv3}, achieved high sample efficiency for policy optimization but were constrained by low-resolution reconstructions that limited their applicability in fine-grained manipulation tasks. Recent advancements have pivoted toward Generative World Models, employing diffusion or transformer backbones to achieve photorealistic simulation. Systems such as Genie~\cite{bruce2024genie} and GameNGen~\cite{valevski2024diffusion} demonstrate the capability to simulate interactive environments or complex game logic with high fidelity. However, a fundamental divergence exists regarding the prediction domain. Joint Embedding Predictive Architectures (JEPA)~\cite{assran2025v,assran2023self} operate in abstract feature spaces to mitigate stochasticity, whereas full generative approaches~\cite{liu2024sora, alonso2024diffusion} synthesize observable data to ensure interpretability. Our framework aligns with the latter strategy but operates within a continuous latent manifold to balance semantic compactness with visual reconstructability.

Despite improved visual fidelity, precise action-controllability and long-term physical consistency remain unresolved challenges. Existing methods generally treat action inputs as passive conditions via embedding concatenation or cross-attention~\cite{zhang2023adding, mou2024t2i}, often resulting in control drift where visual dynamics decouple from input actions over extended horizons. Recent investigations indicate that modulating normalization statistics, as explored in pixel-space diffusion~\cite{yu2025pixeldit}, provides superior structural control compared to standard attention mechanisms. Furthermore, while specialized positional encodings effectively mitigate coordinate drift in discrete autoregressive models~\cite{han2024infinity}, analogous mechanisms for continuous flow fields are currently absent. To address these limitations, we introduce a scale-aware action modulation mechanism for frequency-specific control and formulate a unified spatiotemporal positional encoding to strictly enforce physical consistency during long-term prediction.

